% ============================================================
%  III. PROBLEM FORMULATION
% ============================================================
\section{Problem Formulation}
\label{sec:problem_formulation}

\subsection{System Model}

The \ac{TERESA} platform consists of a Boston Dynamics Spot quadruped and a
Unitree Z1 six-\ac{DOF} serial manipulator arm rigidly mounted on the
robot's back. Four reference frames are defined: the world frame
$\mathcal{F}_W$ (coincident with Spot odometry), the Spot body frame
$\mathcal{F}_B$ (origin at the geometric centre of the base), the arm base
frame $\mathcal{F}_0$ (coincident with Z1 link~0, fixed relative to
$\mathcal{F}_B$), and the end-effector frame $\mathcal{F}_{EE}$ (at the Z1
flange, where the ultrasound probe is attached).

\subsubsection{Spot base kinematics}
The Spot base is modelled as a non-holonomic vehicle constrained to move on a
planar surface. Its configuration is described by
$\boldsymbol{\xi} = [x,\, y,\, \theta]^\top \in \mathbb{R}^2 \times
\mathbb{S}^1$, where $(x, y)$ is the position of $\mathcal{F}_B$ in
$\mathcal{F}_W$ and $\theta$ is the yaw angle. The non-holonomic constraint
restricts the admissible base velocities to
\begin{equation}
  \dot{\boldsymbol{\xi}} =
  \begin{bmatrix} \dot{x} \\ \dot{y} \\ \dot{\theta} \end{bmatrix}
  =
  \begin{bmatrix} \cos\theta \\ \sin\theta \\ 0 \end{bmatrix} v_x
  +
  \begin{bmatrix} 0 \\ 0 \\ 1 \end{bmatrix} \omega_z,
  \label{eq:nonholonomic}
\end{equation}
where $v_x \in \mathbb{R}$ is the forward linear velocity and
$\omega_z \in \mathbb{R}$ is the yaw angular velocity, both expressed in
$\mathcal{F}_B$.

\subsubsection{Body posture}
In addition to planar navigation, Spot can reconfigure its body posture by
changing its height $h \in \mathbb{R}$ and pitch angle
$\varphi \in \mathbb{R}$ relative to the nominal standing configuration.
The height $h$ is a vertical offset of $\mathcal{F}_B$ (negative values lower
the body, bringing the onboard camera closer to the ground); the pitch
$\varphi$ is a rotation about the body $y$-axis (positive values tilt the
body forward, orienting the camera toward the floor). These parameters are
commanded through the native \texttt{body\_pose} interface of the Spot
driver. The combined base configuration is therefore
$[x,\, y,\, \theta,\, h,\, \varphi]^\top$.

\subsubsection{Z1 arm kinematics}
The Z1 arm has $n = 6$ revolute joints with configuration vector
$\mathbf{q} = [q_1, \ldots, q_6]^\top \in \mathbb{R}^6$ and joint velocity
vector $\dot{\mathbf{q}} \in \mathbb{R}^6$. The forward kinematics map
$\mathbf{T}_{EE}(\mathbf{q}) \in SE(3)$ gives the homogeneous transformation
from $\mathcal{F}_0$ to $\mathcal{F}_{EE}$ via the standard product of
exponentials. The arm geometric Jacobian
$\mathbf{J}_{\mathrm{arm}}(\mathbf{q}) \in \mathbb{R}^{6 \times 6}$ relates
joint velocities to the spatial velocity of $\mathcal{F}_{EE}$ expressed in
$\mathcal{F}_0$:
\begin{equation}
  \mathbf{v}_{EE}^{\mathrm{arm}} =
  \mathbf{J}_{\mathrm{arm}}(\mathbf{q})\,\dot{\mathbf{q}}.
  \label{eq:arm_jacobian}
\end{equation}

\subsubsection{Unified control input}
During the approach phase, the system commands the arm joints and the base
velocity simultaneously through the combined input
\begin{equation}
  \mathbf{u} =
  \begin{bmatrix} \dot{\mathbf{q}} \\ v_x \\ \omega_z \end{bmatrix}
  \in \mathbb{R}^8,
  \label{eq:decision_var}
\end{equation}
subject to the velocity bounds
\begin{equation}
  \mathbf{u}_{\min} \leq \mathbf{u} \leq \mathbf{u}_{\max},
  \label{eq:bounds}
\end{equation}
where $\mathbf{u}_{\min}, \mathbf{u}_{\max} \in \mathbb{R}^8$ encode joint
velocity limits for the arm and maximum linear and angular speed for the
base. Body posture commands $(h, \varphi)$ are not part of the continuous
control input; they are set at discrete instants between mission phases,
with the robot given a settle time before the next phase begins.

\subsection{Perception Model}

The quality of patient detection depends on the robot's body posture, as
the onboard Orbbec \ac{RGBD} camera is rigidly attached to $\mathcal{F}_B$.
Let $c(t) \in [0, 1]$ denote the detection confidence provided by the
\ac{YOLO}-based perception pipeline at time $t$. We do not assume an
analytical form for the mapping from posture to confidence, but we rely on
the empirical observation that body height and pitch significantly affect
the camera's field of view and, consequently, the reliability of patient
localisation. The system operates under the premise that there exists a
posture $(h^*, \varphi^*)$ that maximises detection quality for a given
scene geometry, and that this posture can be discovered through active
search.

During scanning, the anatomical target positions are estimated in the
odometry frame $\mathcal{F}_W$ by the anticipatory body scanner
(Section~\ref{sec:active_perception}). These estimates are transformed into
the arm base frame $\mathcal{F}_0$ at execution time via the \ac{TF} tree,
so that changes in body posture $(h, \varphi)$ are automatically reflected
in the \ac{IK} goal without requiring re-estimation of the target.

\subsection{Task Definition}

An autonomous emergency assessment mission requires the \ac{EE} to reach a
sequence of $N$ anatomical target poses. Let
$\mathcal{T} = \{\mathbf{T}_i^* \in SE(3)\}_{i=1}^{N}$ denote the desired
\ac{EE} poses corresponding to the anatomical windows of interest (e.g., the
four \ac{FAST} quadrants). Each target pose
$\mathbf{T}_i^* = (\mathbf{R}_i^*, \mathbf{p}_i^*)$ comprises a desired
\ac{EE} position $\mathbf{p}_i^* \in \mathbb{R}^3$ and orientation
$\mathbf{R}_i^* \in SO(3)$, both expressed in $\mathcal{F}_W$.

The target poses are not known a priori. They are acquired in two stages.
During the approach phase, a wide-field patient detection pipeline provides a
time-varying estimate of the patient's torso position $\hat{\mathbf{p}}(t)$
in the camera frame, which the system transforms to $\mathcal{F}_W$ once the
patient has been localised with sufficient confidence. During the approach
and early scanning phases, an anticipatory multi-view body scanner maps the
patient's anatomy in $\mathcal{F}_W$, from which the target set $\mathcal{T}$
is computed.

\subsection{Problem Statement}

Given the system model defined
in~\eqref{eq:nonholonomic}--\eqref{eq:bounds} and the unknown patient
position, the autonomous emergency assessment problem is to find a control
sequence $\mathbf{u}(t)$ and a sequence of body posture commands
$(h(t), \varphi(t))$ that enable the robot to:
\begin{enumerate}
  \item \textbf{Active perception}: discover a body posture
        $(h^*, \varphi^*, \theta^*)$ from which detection confidence
        exceeds a threshold $c_{\min}$, ensuring robust patient
        localisation before committing to the approach;
  \item \textbf{Navigation}: drive the Spot base from its initial pose to a
        configuration from which all targets in $\mathcal{T}$ lie within
        the reachable workspace of the Z1 arm, using the online patient
        estimate $\hat{\mathbf{p}}(t)$;
  \item \textbf{Body reconfiguration}: for each scanning target $i =
        1,\dots,N$, select a body posture $(h_i, \varphi_i)$ that brings
        the target as close as possible to a preferred workspace centre
        $\mathbf{p}_{\mathrm{sweet}} \in \mathbb{R}^3$ in $\mathcal{F}_0$,
        maximising arm manipulability at execution time;
  \item \textbf{Probe placement}: bring the \ac{EE} to each target pose
        $\mathbf{T}_i^*$ in sequence, satisfying
        \begin{equation}
          \|\mathbf{p}_{EE}(t_i) - \mathbf{p}_i^*\| \leq \varepsilon_p,
          \label{eq:task_constraints}
        \end{equation}
        where $\mathbf{p}_{EE}$ is the current \ac{EE} position and
        $\varepsilon_p > 0$ is the admissible position error;
  \item \textbf{Input feasibility}: satisfy the velocity bounds
        in~\eqref{eq:bounds} and the body posture constraints
        \begin{equation}
          h_{\min} \leq h \leq h_{\max}, \qquad
          0 \leq \varphi \leq \varphi_{\max},
          \label{eq:posture_bounds}
        \end{equation}
        at all times.
\end{enumerate}

\noindent
The following assumptions are adopted throughout this paper:
\begin{itemize}
  \item[\textbf{A1}] The patient remains stationary during the scanning phase.
  \item[\textbf{A2}] Spot operates on flat terrain, so the non-holonomic
        model~\eqref{eq:nonholonomic} is valid.
  \item[\textbf{A3}] The rigid transform between $\mathcal{F}_B$ and
        $\mathcal{F}_0$ is known and constant (fixed mount).
  \item[\textbf{A4}] A single patient is present in the scene.
  \item[\textbf{A5}] The onboard cameras are the only source of perception;
        no external infrastructure (e.g., external tracking systems) is
        available.
\end{itemize}
